\documentclass[a4paper,12pt]{ctexbook}
\usepackage[margin=1cm]{geometry}
\usepackage{graphicx}
\usepackage{subfigure}
\usepackage{float}
\usepackage[colorlinks,linkcolor=black]{hyperref}%colorlinks启用链接颜色，linkcolor指定对应的颜色
\usepackage{listings}
%\usepackage[cache=false]{minted}  % 代码高亮X
\usepackage{fontspec}
\usepackage{tikz,xcolor,mwe}

\definecolor{cvgreen}{HTML}{92D14F}
\definecolor{cvgray}{HTML}{D8E4BE}
\definecolor{cvtext}{HTML}{92909B}
\usetikzlibrary{shadows}

\pagestyle{empty}
\CTEXsetup[format={\large\bfseries}]{section}%可以让section的标题左对齐。
\CTEXsetup[format={\Large\bfseries}]{chapter}%可以让chapter的标题左对齐。
%\CTEXsetup[format+={\flushleft}]{section}%让section的标题居左
%\renewcommand{\thesection}{\chinese{section}}%将“1.1”改为汉字“一”，但是subsection就会变成  六.1 ，比较难看，还是不用比较好。

\begin{document}


\begin{tikzpicture}[remember picture,overlay]
\fill[cvgreen] (current page.north west) rectangle ([xshift=6cm]current page.south west);% green bar
\fill[cvgray] ([yshift=-11cm]current page.north west) rectangle ([yshift=-17cm]current page.north east); % gray bar

\node[cvtext,right] at ([xshift=3.5cm,yshift=-13cm]current page.north west) {\Huge My ML way};
\node[cvtext,above left] at ([xshift=-1cm,yshift=-16.5cm]current page.north east) {\Large\bfseries \today};
% cover photo
\node[inner sep=0pt,below right] (image) at ([xshift=17cm,yshift=-1cm]current page.north west) {\includegraphics[width=3cm]{hainulogo.png}};
% name and address
\node[fill=white,align=center,text width=6.4cm,inner sep=0.8cm,below] (name) at (image.south) {};
\node[text width=15cm,inner sep=0.3cm,below right] at (name.south west){\Large Author:Flynn\\E-Mail:zofon@qq.com};
% attachments
\node[white,text width=5cm,inner sep=0.6cm,above right] at ([yshift=1cm]current page.south west)
{\large\obeylines\textbf{This is\\Open Source}};
\end{tikzpicture}

\tableofcontents

\chapter{相关的概念}
在学习机器学习之前首先要了解一些相关的概念。
\begin{description}
  \item[AI] 人工智能是最早出现的，也是最大、最外侧的同心圆；
  \item[ML] 机器学习
  \item[DL] 深度学习，深度学习造成了前所未有的巨大的影响， 是当今人工智能大爆炸的核心驱动。
\end{description}
五十年代，人工智能曾一度被极为看好。之后，人工智能的一些较小的子集发展了起来。先是机器学习，然后是深度学习。深度学习又是机器学习的子集。。

早在1956年夏天那次会议，人工智能的先驱们就梦想着用当时刚刚出现的计算机来构造复杂的、拥有与人类智慧同样本质特性的机器。这就是我们现在所说的“强人工智能”（General AI）。这个无所不能的机器，它有着我们所有的感知（甚至比人更多），我们所有的理性，可以像我们一样思考。

人们在电影里也总是看到这样的机器：友好的，像星球大战中的C-3PO；邪恶的，如终结者。强人工智能现在还只存在于电影和科幻小说中，原因不难理解，我们还没法实现它们，至少目前还不行。

我们目前能实现的，一般被称为“弱人工智能”（Narrow AI）。弱人工智能是能够与人一样，甚至比人更好地执行特定任务的技术。例如，Pinterest上的图像分类；或者Facebook的人脸识别。

这些是弱人工智能在实践中的例子。这些技术实现的是人类智能的一些具体的局部。但它们是如何实现的？这种智能是从何而来？这就带我们来到同心圆的里面一层，机器学习。

\input{chapter\_如何研究一个机器学习算法.tex}

\input{chapter\_框架比较.tex}

\input{chapter\_神经网络入门.tex}

\input{chapter\_深度学习实例.tex}

\bibliographystyle{plain}
\bibliography{Referenzarchiv}
\end{document} 